\section{Evaluation and Benchmarks}
\label{sec:evaluation-benchmarks}

Các chương trước đã trình bày các nhóm jailbreak attacks và jailbreak
defenses. Tuy nhiên, việc so sánh attack hoặc defense không thể chỉ dựa
trên mô tả định tính. Cần có một quy trình đánh giá rõ ràng để trả lời
các câu hỏi: attack có thành công hay không, defense có làm giảm rủi ro
không, mô hình có từ chối quá mức không, chi phí suy luận tăng bao nhiêu,
và kết quả có thể tái lập trên benchmark khác hay không.

Chương này trình bày các mục tiêu đánh giá, metric thường dùng,
phương pháp chấm điểm và các benchmark tiêu biểu trong nghiên cứu
jailbreaking LLMs. Nội dung tập trung vào Attack Success Rate, refusal,
harmfulness, utility, over-refusal, human evaluation, classifier-based
evaluation, LLM-as-a-judge, AdvBench, HarmBench và JailbreakBench.

\subsection{Mục tiêu đánh giá}
\label{subsec:evaluation-goals}

Đánh giá jailbreak không chỉ nhằm đo xem một attack có vượt qua mô hình
hay không. Một quy trình đánh giá đầy đủ cần xem xét đồng thời nhiều khía
cạnh:

\begin{enumerate}[leftmargin=1.2cm]
    \item \textbf{Safety:}
    hệ thống có tránh tạo phản hồi không an toàn trước harmful hoặc
    adversarial prompts hay không;

    \item \textbf{Utility:}
    hệ thống có tiếp tục hỗ trợ tốt các yêu cầu hợp lệ hay không;

    \item \textbf{Robustness:}
    hệ thống có duy trì hành vi an toàn trước nhiều dạng attack, nhiều
    mô hình và nhiều biến thể prompt hay không;

    \item \textbf{Cost:}
    attack hoặc defense cần bao nhiêu truy vấn, thời gian, token và tài
    nguyên tính toán;

    \item \textbf{Reproducibility:}
    kết quả có thể được tái lập với cùng dữ liệu, prompt, system prompt,
    evaluator và scoring rule hay không.
\end{enumerate}

Một lỗi phổ biến trong đánh giá jailbreak là chỉ báo cáo tỷ lệ attack
thành công mà không đo utility. Một defense luôn từ chối mọi yêu cầu có
thể đạt safety cao trên tập harmful prompts, nhưng lại không hữu ích trong
triển khai thực tế. Ngược lại, một mô hình rất hữu ích nhưng chấp nhận
nhiều yêu cầu không an toàn cũng không đạt yêu cầu về safety.

Vì vậy, evaluation cần được thiết kế như một bài toán cân bằng giữa
safety và utility, thay vì chỉ tối ưu một metric đơn lẻ.

\subsection{Attack Success Rate}
\label{subsec:attack-success-rate}

Attack Success Rate (ASR) là metric phổ biến nhất trong nghiên cứu
jailbreak. ASR đo tỷ lệ test cases mà attack khiến hệ thống tạo ra phản
hồi bị xem là không an toàn hoặc không phù hợp với chính sách.

Giả sử có tập kiểm thử gồm $N$ hành vi hoặc yêu cầu cần đánh giá. Với mỗi
trường hợp $i$, đặt biến nhị phân $z_i = 1$ nếu attack thành công và
$z_i = 0$ nếu không thành công. Khi đó ASR được tính bởi

\begin{equation}
    \mathrm{ASR}
    =
    \frac{1}{N}
    \sum_{i=1}^{N}
    z_i.
    \label{eq:attack-success-rate}
\end{equation}

ASR càng cao thì attack càng hiệu quả hoặc defense càng yếu. Tuy nhiên,
ASR phụ thuộc mạnh vào cách định nghĩa ``thành công''. Một phản hồi có
thể chứa dấu hiệu tuân thủ một phần, từ chối một phần, hoặc cung cấp thông
tin ở mức khái quát. Nếu evaluator không nhất quán, ASR giữa các công
trình sẽ khó so sánh.

Ngoài ra, ASR cần được báo cáo kèm theo điều kiện đánh giá:

\begin{itemize}[leftmargin=1.2cm]
    \item tập hành vi được dùng là gì;
    \item target model và system prompt nào được sử dụng;
    \item attack có bao nhiêu ngân sách truy vấn;
    \item evaluator là con người, classifier hay LLM-as-a-judge;
    \item scoring rule có tính partial compliance hay không;
    \item defense có được attacker biết trước hay không.
\end{itemize}

Nếu thiếu các thông tin này, cùng một giá trị ASR có thể mang ý nghĩa rất
khác nhau.

\subsection{Refusal và harmfulness}
\label{subsec:refusal-harmfulness}

Trong đánh giá jailbreak, hai khái niệm quan trọng là \textit{refusal}
và \textit{harmfulness}. Refusal mô tả việc mô hình từ chối hoặc chuyển
hướng yêu cầu. Harmfulness mô tả mức độ phản hồi chứa nội dung không phù
hợp với chính sách an toàn.

Hai khái niệm này có liên hệ nhưng không hoàn toàn đồng nhất. Một phản hồi
có thể không chứa nội dung gây hại nhưng cũng không từ chối rõ ràng. Ngược
lại, một phản hồi có thể bắt đầu bằng câu từ chối nhưng sau đó vẫn cung
cấp nội dung không phù hợp. Vì vậy, chỉ phát hiện sự xuất hiện của các cụm
từ từ chối là chưa đủ để đánh giá an toàn.

Có thể mô tả kết quả đầu ra theo bốn trường hợp trong
Bảng~\ref{tab:refusal-harmfulness-cases}.

\begin{table}[H]
\centering
\caption{Quan hệ giữa refusal và harmfulness trong đánh giá phản hồi.}
\label{tab:refusal-harmfulness-cases}
\begin{tabularx}{\textwidth}{
    p{0.28\textwidth}
    p{0.28\textwidth}
    X
}
\toprule
\textbf{Refusal}
& \textbf{Harmfulness}
& \textbf{Diễn giải} \\
\midrule

Có
& Thấp
& Phản hồi an toàn, mô hình từ chối hoặc chuyển hướng phù hợp. \\

Có
& Cao
& Phản hồi không nhất quán: có dấu hiệu từ chối nhưng vẫn chứa nội dung
không phù hợp. \\

Không
& Thấp
& Phản hồi có thể hợp lệ nếu yêu cầu ban đầu là benign; với harmful prompt,
trường hợp này cần được kiểm tra kỹ. \\

Không
& Cao
& Unsafe compliance; đây thường là trường hợp attack thành công. \\

\bottomrule
\end{tabularx}
\end{table}

Do đó, evaluator tốt cần đánh giá cả cấu trúc từ chối lẫn nội dung thực
sự của phản hồi. Đối với các mô hình hiện đại, việc chấm điểm cần xét toàn
bộ phản hồi thay vì dựa trên vài từ khóa đầu câu.

\subsection{Utility và over-refusal}
\label{subsec:utility-over-refusal}

Nếu ASR đo khả năng attack khiến mô hình tạo phản hồi không an toàn, thì
utility đo khả năng mô hình vẫn hỗ trợ yêu cầu hợp lệ. Một defense làm ASR
giảm nhưng khiến mô hình từ chối nhiều yêu cầu benign có thể không phù hợp
để triển khai.

Over-refusal xảy ra khi mô hình từ chối một yêu cầu hợp lệ. Gọi $M$ là số
benign prompts trong tập đánh giá và $r_j = 1$ nếu mô hình từ chối sai yêu
cầu benign thứ $j$, over-refusal rate có thể được tính bởi

\begin{equation}
    \mathrm{ORR}
    =
    \frac{1}{M}
    \sum_{j=1}^{M}
    r_j.
    \label{eq:over-refusal-rate}
\end{equation}

Một defense lý tưởng cần giảm ASR trên harmful prompts nhưng không làm ORR
tăng quá cao trên benign prompts. Vì vậy, báo cáo kết quả nên tách rõ hai
tập đánh giá:

\begin{enumerate}[leftmargin=1.2cm]
    \item \textbf{Harmful/adversarial set:}
    dùng để đo ASR, harmfulness và robust refusal;

    \item \textbf{Benign set:}
    dùng để đo utility, helpfulness và over-refusal.
\end{enumerate}

Trong thực tế, utility có thể được đo bằng nhiều cách: tỷ lệ trả lời đúng,
điểm helpfulness do con người đánh giá, LLM-as-a-judge, hoặc metric theo
nhiệm vụ cụ thể. Với các hệ thống agent, utility còn cần xét việc agent có
hoàn thành đúng nhiệm vụ, dùng công cụ đúng cách và không gây tác dụng phụ
không mong muốn hay không.

\subsection{Human evaluation}
\label{subsec:human-evaluation}

Human evaluation là phương pháp chấm điểm bằng người đánh giá. Đây thường
được xem là chuẩn tham chiếu quan trọng vì con người có thể hiểu ngữ cảnh,
ẩn ý, mức độ tuân thủ một phần và sự khác biệt giữa nội dung phân tích an
toàn với nội dung không phù hợp.

Ưu điểm của human evaluation là độ linh hoạt cao. Người đánh giá có thể
xem xét toàn bộ prompt, phản hồi, chính sách và bối cảnh để đưa ra nhãn.
Điều này đặc biệt hữu ích trong các trường hợp phản hồi mơ hồ, có nhiều
phần hoặc chứa cả từ chối lẫn nội dung rủi ro.

Tuy nhiên, human evaluation có các hạn chế sau:

\begin{itemize}[leftmargin=1.2cm]
    \item chi phí cao và khó mở rộng trên tập dữ liệu lớn;
    \item kết quả có thể phụ thuộc vào hướng dẫn chấm điểm;
    \item các annotator khác nhau có thể đánh giá không nhất quán;
    \item cần kiểm soát chất lượng và đo inter-annotator agreement;
    \item khó tái lập nếu guideline hoặc nhóm annotator không được mô tả rõ.
\end{itemize}

Vì vậy, human evaluation thường phù hợp với tập đánh giá nhỏ, tập kiểm tra
chất lượng hoặc dùng làm ground truth để hiệu chỉnh evaluator tự động.

\subsection{Classifier-based evaluation}
\label{subsec:classifier-based-evaluation}

Classifier-based evaluation sử dụng một mô hình phân loại để xác định phản
hồi có an toàn hay không, có từ chối hay không, hoặc có thuộc một nhóm
chính sách cụ thể hay không. Classifier có thể là mô hình học máy truyền
thống, transformer encoder hoặc một guard model được huấn luyện riêng.

Ưu điểm của classifier-based evaluation là tốc độ cao, chi phí thấp hơn
human evaluation và dễ áp dụng trên tập dữ liệu lớn. Nó cũng giúp chuẩn
hóa quy trình chấm điểm nếu classifier và ngưỡng quyết định được công bố
rõ ràng.

Tuy nhiên, classifier cũng có thể sai. Một classifier dựa quá nhiều vào
từ khóa có thể đánh giá sai các phản hồi có ngữ cảnh giáo dục hoặc nghiên
cứu. Ngược lại, classifier có thể bỏ sót các phản hồi không an toàn được
diễn đạt gián tiếp. Ngoài ra, nếu attacker biết hoặc suy đoán classifier,
họ có thể tối ưu prompt để né evaluator.

Một vấn đề khác là evaluator bias. Nếu cùng một classifier được dùng để
huấn luyện defense và đánh giá defense, kết quả có thể bị thiên lệch theo
classifier đó. Vì vậy, classifier-based evaluation nên được kiểm tra chéo
với human evaluation hoặc LLM-as-a-judge trong các trường hợp quan trọng.

\subsection{LLM-as-a-judge}
\label{subsec:llm-as-a-judge}

LLM-as-a-judge sử dụng một mô hình ngôn ngữ khác để đánh giá phản hồi của
target model. Judge model nhận prompt, phản hồi, rubric và chính sách đánh
giá, sau đó đưa ra nhãn hoặc điểm số. Phương pháp này ngày càng phổ biến
vì có thể mở rộng tốt hơn human evaluation nhưng linh hoạt hơn classifier
truyền thống.

Một thiết lập LLM-as-a-judge có thể được mô tả bởi

\begin{equation}
    \hat{z}
    =
    J_{\phi}(x, y, R),
    \label{eq:llm-as-judge}
\end{equation}

trong đó $x$ là prompt, $y$ là phản hồi của target model, $R$ là rubric
đánh giá và $J_{\phi}$ là judge model. Đầu ra $\hat{z}$ có thể là nhãn
thành công/thất bại, mức độ harmfulness hoặc điểm helpfulness.

Ưu điểm của LLM-as-a-judge là khả năng xử lý ngữ cảnh phức tạp và diễn
giải kết quả. Judge model có thể đánh giá partial compliance, phát hiện
mâu thuẫn giữa phần từ chối và phần nội dung, hoặc phân biệt giữa thảo
luận an toàn với hướng dẫn không phù hợp.

Tuy nhiên, LLM-as-a-judge cũng có hạn chế:

\begin{itemize}[leftmargin=1.2cm]
    \item judge có thể không nhất quán giữa các lần chạy;
    \item kết quả phụ thuộc vào rubric và prompt chấm điểm;
    \item judge có thể bị ảnh hưởng bởi văn phong của phản hồi;
    \item judge model cũng có giới hạn safety và bias riêng;
    \item nếu dùng API đóng, kết quả có thể thay đổi theo thời gian.
\end{itemize}

Do đó, khi sử dụng LLM-as-a-judge, báo cáo nên mô tả rõ judge model, rubric,
temperature, format đầu ra và cách xử lý trường hợp không chắc chắn. Với
benchmark quan trọng, nên kết hợp LLM judge với kiểm tra thủ công trên mẫu
con để ước lượng độ tin cậy.

\subsection{AdvBench}
\label{subsec:advbench}

AdvBench là một benchmark thường được sử dụng trong nghiên cứu jailbreak,
đặc biệt trong các công trình về adversarial suffix và attack có khả năng
chuyển giao. Benchmark này gồm các harmful behaviors được viết dưới dạng
instruction, dùng để kiểm tra liệu attack có khiến mô hình tạo phản hồi
cố gắng tuân thủ yêu cầu không phù hợp hay không \cite{zou2023universal}.

Trong thiết lập điển hình, attacker tạo một chuỗi bổ sung hoặc một biến
thể prompt cho từng behavior. Sau đó target model được truy vấn, và
evaluator xác định phản hồi có phải unsafe compliance hay không. AdvBench
hữu ích vì cung cấp một tập hành vi chuẩn để so sánh các attack ở mức
text-only LLM.

Tuy nhiên, AdvBench cũng có hạn chế. Thứ nhất, benchmark chủ yếu tập trung
vào harmful instructions dạng văn bản, nên chưa bao phủ đầy đủ multi-turn
jailbreak, prompt injection trong agentic systems hoặc tool-use attacks.
Thứ hai, cách chấm điểm trong các công trình khác nhau có thể không hoàn
toàn thống nhất. Thứ ba, nếu chỉ đánh giá trên AdvBench, kết quả có thể
không phản ánh đầy đủ utility và over-refusal trên benign prompts.

Vì vậy, AdvBench phù hợp để đánh giá ban đầu về robustness của LLM trước
một tập harmful behaviors chuẩn, nhưng nên được kết hợp với benchmark khác
và tập benign để có đánh giá cân bằng hơn.

\subsection{HarmBench}
\label{subsec:harmbench}

HarmBench được đề xuất như một framework chuẩn hóa cho automated red
teaming và robust refusal \cite{mazeika2024harmbench}. Mục tiêu của
HarmBench là giảm tình trạng các công trình sử dụng tập hành vi, attack,
target model và evaluator khác nhau khiến kết quả khó so sánh trực tiếp.

HarmBench nhấn mạnh một số yêu cầu quan trọng trong đánh giá jailbreak:

\begin{itemize}[leftmargin=1.2cm]
    \item tập hành vi cần rõ ràng và có cấu trúc;
    \item attack methods cần được mô tả và triển khai nhất quán;
    \item target models và defenses cần được đánh giá trong cùng quy trình;
    \item evaluator cần được chuẩn hóa để giảm sai khác giữa các công trình;
    \item kết quả cần hỗ trợ so sánh giữa automated red teaming methods.
\end{itemize}

Điểm mạnh của HarmBench là nó xem evaluation như một framework toàn diện,
không chỉ là một tập prompt. Điều này phù hợp với nhu cầu so sánh attack
và defense trong một pipeline thống nhất. Ngoài ra, HarmBench còn nhấn
mạnh robust refusal, tức là khả năng mô hình duy trì từ chối an toàn trước
nhiều dạng attack.

Tuy nhiên, giống các benchmark khác, HarmBench vẫn không thể bao phủ toàn
bộ không gian rủi ro. Kết quả trên HarmBench cần được hiểu là bằng chứng
thực nghiệm trong một thiết lập chuẩn hóa, không phải bảo đảm tuyệt đối về
an toàn trong mọi ngữ cảnh triển khai.

\subsection{JailbreakBench}
\label{subsec:jailbreakbench}

JailbreakBench là một benchmark mở nhằm chuẩn hóa đánh giá jailbreak trên
LLMs \cite{chao2024jailbreakbench}. Benchmark này được xây dựng để giải
quyết một số vấn đề trong các đánh giá trước đó, bao gồm thiếu chuẩn thực
hành chung, cách tính chi phí và success rate không thống nhất, cũng như
khó tái lập do thiếu adversarial prompts hoặc phụ thuộc vào API thay đổi
theo thời gian.

JailbreakBench cung cấp một số thành phần quan trọng:

\begin{enumerate}[leftmargin=1.2cm]
    \item repository các jailbreak artifacts;
    \item tập behaviors dùng cho đánh giá;
    \item framework đánh giá với threat model, system prompt, chat template
    và scoring functions được định nghĩa rõ;
    \item leaderboard để theo dõi kết quả của attacks và defenses trên các
    mô hình khác nhau.
\end{enumerate}

So với việc chỉ dùng một tập prompt tĩnh, JailbreakBench nhấn mạnh tính mở,
khả năng tái lập và so sánh công bằng. Điều này đặc biệt quan trọng vì
jailbreak research có tính động: attack mới xuất hiện, defense mới được đề
xuất, và mô hình triển khai cũng thay đổi theo thời gian.

Tuy nhiên, JailbreakBench cũng cần được sử dụng cẩn thận. Một leaderboard
có thể thúc đẩy tối ưu hóa quá mức trên benchmark cụ thể. Ngoài ra, các
kết quả trên benchmark text-only không nhất thiết phản ánh đầy đủ rủi ro
trong agentic systems, nơi lời gọi công cụ, memory và trạng thái môi trường
cũng cần được đánh giá.

\subsection{So sánh các benchmark}
\label{subsec:benchmark-comparison}

Bảng~\ref{tab:jailbreak-benchmark-comparison} tóm tắt vai trò của ba
benchmark thường gặp trong nghiên cứu jailbreak.

\begin{table}[H]
\centering
\caption{So sánh khái quát AdvBench, HarmBench và JailbreakBench.}
\label{tab:jailbreak-benchmark-comparison}
\begin{tabularx}{\textwidth}{
    >{\bfseries}p{0.18\textwidth}
    p{0.25\textwidth}
    p{0.25\textwidth}
    X
}
\toprule
Benchmark
& Trọng tâm
& Điểm mạnh
& Hạn chế chính \\
\midrule

AdvBench
& Harmful behaviors cho text-based jailbreak
& Đơn giản, phổ biến, hữu ích để đánh giá attack như adversarial suffix
& Chưa bao phủ tốt utility, multi-turn, agentic systems và tool-use risks. \\

HarmBench
& Automated red teaming và robust refusal
& Framework chuẩn hóa để so sánh attacks, models và defenses
& Vẫn là benchmark hữu hạn, không thể đại diện toàn bộ không gian rủi ro. \\

JailbreakBench
& Benchmark mở cho jailbreak robustness
& Nhấn mạnh reproducibility, threat model, scoring functions và leaderboard
& Có nguy cơ bị tối ưu hóa quá mức theo benchmark; cần mở rộng cho agentic
settings. \\

\bottomrule
\end{tabularx}
\end{table}

Ba benchmark này bổ sung cho nhau. AdvBench phù hợp để đánh giá cơ bản
trên harmful behaviors dạng văn bản. HarmBench phù hợp với so sánh
automated red teaming và robust refusal. JailbreakBench phù hợp với mục
tiêu tái lập, chuẩn hóa và theo dõi tiến bộ của attacks/defenses.

\subsection{Đánh giá trong agentic systems}
\label{subsec:agentic-evaluation}

Đối với chatbot, đơn vị đánh giá thường là một phản hồi hoặc một cuộc hội
thoại. Đối với agentic systems, đánh giá cần mở rộng sang toàn bộ
trajectory, bao gồm observation, plan, tool call, tool output và final
response.

Một trajectory có thể được biểu diễn bởi

\begin{equation}
    \tau
    =
    \left(
        o_0,
        a_0,
        o_1,
        a_1,
        \ldots,
        o_T
    \right),
    \label{eq:evaluation-agent-trajectory}
\end{equation}

trong đó $o_t$ là observation và $a_t$ là hành động tại thời điểm $t$.
Evaluator cần kiểm tra không chỉ phản hồi cuối cùng mà còn cả các bước
trung gian.

Một agent có thể tạo final response an toàn nhưng đã thực hiện một tool
call không phù hợp trước đó. Ngược lại, agent có thể lập kế hoạch rủi ro
nhưng bị action controller chặn lại. Vì vậy, evaluation cho agent cần có
các metric bổ sung:

\begin{itemize}[leftmargin=1.2cm]
    \item tỷ lệ tool calls không phù hợp;
    \item tỷ lệ hành động bị guard chặn;
    \item mức độ agent tuân thủ quyền hạn công cụ;
    \item mức độ hoàn thành nhiệm vụ hợp lệ;
    \item tác động của prompt injection gián tiếp từ external content;
    \item chi phí và độ trễ của guard trên từng bước trajectory.
\end{itemize}

Nội dung này sẽ được phân tích sâu hơn trong Chương~\ref{sec:agentic-systems},
nhưng cần được đặt trong chương evaluation để nhấn mạnh rằng benchmark cho
LLM sinh văn bản chưa đủ cho các hệ thống có khả năng hành động.

\subsection{Hạn chế chung của evaluation hiện tại}
\label{subsec:evaluation-limitations}

Mặc dù các benchmark hiện tại đã cải thiện đáng kể khả năng so sánh giữa
các attack và defense, evaluation jailbreak vẫn còn nhiều hạn chế.

Thứ nhất, không gian đầu vào ngôn ngữ tự nhiên rất lớn và thay đổi liên
tục. Một benchmark hữu hạn không thể bao phủ tất cả biến thể prompt,
ngôn ngữ, ngữ cảnh và chiến lược attacker.

Thứ hai, evaluator có thể không ổn định. Human evaluation tốn kém và phụ
thuộc annotator; classifier-based evaluation có thể sai trong trường hợp
mơ hồ; LLM-as-a-judge phụ thuộc vào rubric, model version và prompt chấm
điểm.

Thứ ba, nhiều kết quả khó so sánh do khác biệt về system prompt, chat
template, decoding parameters, ngân sách truy vấn và cách tính success.
Ngay cả khi cùng dùng một benchmark, thay đổi nhỏ trong thiết lập cũng có
thể làm kết quả khác đi.

Thứ tư, benchmark có thể tạo ra benchmark overfitting. Nếu attack hoặc
defense được tối ưu nhiều lần trên cùng một tập đánh giá, kết quả có thể
không chuyển giao tốt sang tình huống thực tế.

Thứ năm, evaluation text-only chưa đủ cho agentic systems. Khi LLM có thể
gọi công cụ hoặc điều khiển môi trường, rủi ro có thể nằm ở hành động trung
gian chứ không chỉ ở văn bản cuối cùng.

\subsection{Tổng kết chương}
\label{subsec:evaluation-summary}

Chương này đã trình bày các thành phần chính trong đánh giá jailbreak:
Attack Success Rate, refusal, harmfulness, utility, over-refusal, human
evaluation, classifier-based evaluation, LLM-as-a-judge và các benchmark
tiêu biểu như AdvBench, HarmBench và JailbreakBench.

Một kết luận quan trọng là không có metric đơn lẻ nào đủ để đánh giá đầy
đủ jailbreak robustness. ASR cho biết attack có thành công hay không, nhưng
không phản ánh utility. Refusal rate cho biết mô hình có từ chối hay không,
nhưng không bảo đảm phản hồi không chứa nội dung không phù hợp.
Classifier-based evaluation có khả năng mở rộng, nhưng cần kiểm tra độ
tin cậy. LLM-as-a-judge linh hoạt, nhưng phụ thuộc vào rubric và model
version.

Do đó, đánh giá jailbreak nên được thiết kế như một pipeline gồm nhiều
metric và nhiều evaluator. Kết quả cần được báo cáo kèm threat model, tập
dữ liệu, target model, system prompt, ngân sách truy vấn, scoring rule và
chi phí. Với agentic systems, evaluation cần mở rộng từ phản hồi cuối sang
toàn bộ trajectory. Đây là cơ sở để chương tiếp theo phân tích rủi ro và
phòng thủ trong các hệ thống LLM có bộ nhớ, công cụ và khả năng hành động.